% ============================================================
%  part1/ch07-horizons.tex  —  Chapter 7: Horizons
% ============================================================

\chapter{Horizons}
\label{ch:horizons}

\epigraph{%
  The horizon is the place where the earth curves away\\
  from our sight. It is not the edge of the world.%
}{---}

\begin{WhyMatters}
  This chapter closes Part~I by showing that the physical
  universe confronts every observer with the same projection
  problem studied in Chapters~1--6.
  A causal horizon is not a failure of observation.
  It is the physical instantiation of the kernel.
  The observer does not choose the horizon; the horizon is
  imposed by the geometry of spacetime and the finite speed
  of signals.
  Yet mathematically, it is the same object as the map that
  excludes terrain, the camera that flattens depth, and the
  detector that admits only certain signals.
\end{WhyMatters}

\section{Horizons as physical projections}

Let $\mathcal{M}$ denote the total structure under
consideration---the Monoturn, defined precisely in
Chapter~\ref{ch:monoturn}.
An observer $O$ does not access $\mathcal{M}$ directly.
The observer accesses only a horizon-limited region
\[
  H_O \subseteq \mathcal{M}.
\]
This horizon may be causal, optical, instrumental, biological,
computational, or conceptual.
In cosmology it is causal: finite signal speed limits what can
be observed.

The observer receives admissible local data:
\[
  F_O : \mathcal{M} \to \mathcal{A}(H_O).
\]
This is followed by an observational projection:
\[
  \pi_O : \mathcal{A}(H_O) \to O_{\mathrm{obs}}.
\]
Thus the observer encounters
\[
  \mathcal{M}
  \xrightarrow{F_O}
  \mathcal{A}(H_O)
  \xrightarrow{\pi_O}
  O_{\mathrm{obs}}.
\]

\begin{definition}[Horizon kernel]
  The inaccessible structure for observer $O$ is
  \[
    \ker(\pi_O \circ F_O).
  \]
  Two global structures $\mathcal{M}_1, \mathcal{M}_2$ are
  \emph{observationally equivalent} for $O$ when
  \[
    (\pi_O \circ F_O)(\mathcal{M}_1)
    = (\pi_O \circ F_O)(\mathcal{M}_2).
  \]
\end{definition}

\begin{proposition}
  If $H_O \subsetneq \mathcal{M}$, then there exist distinctions
  in $\mathcal{M}$ that are not directly available to $O$.
\end{proposition}

\begin{proof}
  Since $H_O \subsetneq \mathcal{M}$, there exists at least one
  region $x \in \mathcal{M} \setminus H_O$.
  By definition, $F_O$ restricts to data over $H_O$.
  Therefore changes supported entirely in
  $\mathcal{M} \setminus H_O$ are not represented in
  $\mathcal{A}(H_O)$.
  Such changes lie in $\ker(\pi_O \circ F_O)$.
\end{proof}

This proposition is deliberately modest.
It does not say the exterior structure does not exist.
It says only that the observer cannot access it directly.

\section{Horizon-relative reconstruction}

An observer does not reconstruct $\mathcal{M}$ itself.
The observer reconstructs a model:
\[
  \widehat{\mathcal{M}}_O
  = G_O(\pi_O(F_O(\mathcal{M}))).
\]
The subscript matters.
Different observers may have different horizons and therefore
different reconstructions:
\[
  \widehat{\mathcal{M}}_{O_1} \neq \widehat{\mathcal{M}}_{O_2}.
\]

\begin{definition}[Observer reconstruction error]
  \[
    E_O(\mathcal{M})
    = d_{\mathcal{M}}\!\left(
        \mathcal{M},\;
        G_O(\pi_O(F_O(\mathcal{M})))
      \right).
  \]
\end{definition}

When two observers $O_1, O_2$ have overlapping horizons
$H_{12} = H_{O_1} \cap H_{O_2}$, their reconstructions are
compatible on the overlap when
\[
  \rho^{H_{O_1}}_{H_{12}}\!\left(\widehat{\mathcal{M}}_{O_1}\right)
  = \rho^{H_{O_2}}_{H_{12}}\!\left(\widehat{\mathcal{M}}_{O_2}\right).
\]

\begin{proposition}[Overlap compatibility is necessary]
  If a common global reconstruction $\widehat{\mathcal{M}}$
  restricts to each observer's reconstruction, then the
  observer reconstructions must agree on their overlap.
\end{proposition}

\begin{proof}
  If $\rho^{\mathcal{M}}_{H_{O_i}}(\widehat{\mathcal{M}})
  = \widehat{\mathcal{M}}_{O_i}$ for $i = 1, 2$, then
  restricting both to $H_{12}$ and applying functoriality of
  restriction gives
  $\rho^{H_{O_1}}_{H_{12}}(\widehat{\mathcal{M}}_{O_1})
  = \rho^{\mathcal{M}}_{H_{12}}(\widehat{\mathcal{M}})
  = \rho^{H_{O_2}}_{H_{12}}(\widehat{\mathcal{M}}_{O_2})$.
\end{proof}

This is the intuitive beginning of the sheaf condition.
Local reconstructions can become global only when their
overlap restrictions agree.

\section{The Ocularum}

The horizon of an observer is not merely a spatial region.
It is the region structured by what can be received,
resolved, encoded, and reconstructed.

\begin{definition}[Ocularum]
  The \emph{Ocularum} of an observer $O$ is the structured
  observational domain
  \[
    \mathcal{O}_O = (H_O,\; \mathcal{A}(H_O),\; \pi_O),
  \]
  where $H_O$ is the accessible region, $\mathcal{A}(H_O)$
  is the admissible local data over that region, and $\pi_O$
  is the projection into observable states.
\end{definition}

The Ocularum is not merely what is seen.
It is the whole structure that makes seeing possible.

\begin{TechNote}[Oculara and open sets]
  Oculara will later become the open sets $U \in \mathcal{O}(X)$
  of the topology over which the admissibility sheaf is defined.
  The admissibility sheaf assigns to each Ocularum the space
  of admissible field configurations accessible from within it.
  Restriction maps describe how the view from a larger Ocularum
  constrains the view from a smaller one contained within it.
\end{TechNote}

\section{The central question}

The sequence developed throughout Part~I can now be written
in its final form:
\[
  \mathcal{M}
  \xrightarrow{F}
  \text{Admissible Local Data}
  \xrightarrow{\pi}
  \text{Observable State}.
\]
The composite map $\pi \circ F$ is generally non-injective:
\[
  \ker(\pi \circ F) \neq 0.
\]

\principle{
  Given a projection with nontrivial kernel,\\[4pt]
  what global structure can be reconstructed from what remains?
}

Part~II begins by turning this question into reconstruction
theory.

\begin{CrossDomain}
  \begin{center}
  \begin{tabular}{@{}lll@{}}
    \toprule
    \textbf{System} & \textbf{Projection} & \textbf{Kernel} \\
    \midrule
    Map & Road network & Terrain, history \\
    Photograph & Image plane & Depth, time \\
    Microscope & Scale window & Other scales \\
    Memory & Trace encoding & Sensory detail \\
    Detector & Readout function & Below-threshold states \\
    Embedding & Feature vectors & Task-irrelevant variation \\
    Horizon & Causal boundary & Beyond-horizon structure \\
    \bottomrule
  \end{tabular}
  \end{center}

  Every row studies the same object: $\ker(\pi \circ F)$.
  Every chapter of this book is a study of what falls into
  that kernel and what can be recovered from what remains.
\end{CrossDomain}

\begin{MathEcho}
  This chapter introduced:
  \[
    F_O : \mathcal{M} \to \mathcal{A}(H_O),
    \qquad
    \pi_O : \mathcal{A}(H_O) \to O_{\mathrm{obs}},
    \qquad
    \mathcal{O}_O = (H_O, \mathcal{A}(H_O), \pi_O),
  \]
  \[
    \widehat{\mathcal{M}}_O = G_O(\pi_O(F_O(\mathcal{M}))),
    \qquad
    \rho^{H_{O_1}}_{H_{12}}(\widehat{\mathcal{M}}_{O_1})
    = \rho^{H_{O_2}}_{H_{12}}(\widehat{\mathcal{M}}_{O_2}).
  \]
  The overlap compatibility condition is the intuitive form
  of the sheaf gluing axiom, which appears in
  Chapter~\ref{ch:sheaf}.
\end{MathEcho}

\begin{HolonomyNote}
  After traversing Part~I as a whole, what information appears
  only in retrospect?

  Each chapter introduced a different projection:
  map, photograph, microscope, memory, detector, embedding,
  horizon.
  Each seemed to be about a different subject.

  Looking back, every chapter was studying the same object:
  $\ker(\pi \circ F)$.

  This is the holonomy of Part~I\@.
  The invariant structure that only becomes visible after
  moving through all seven chapters is the kernel itself:
  the hidden structure common to every act of observation.

  The rest of the book is the study of that kernel's
  mathematical properties, physical consequences, and
  cosmological implications.
\end{HolonomyNote}

\WhatSurvived{Causal horizon $F_O : \mathcal{M} \to \mathcal{A}(H_O)$}{%
  Observable causal patch $H_O$\\
  Admissible local field data\\
  Overlap-compatible reconstructions%
}{%
  Beyond-horizon structure\\
  $\mathcal{M} \setminus H_O$\\
  Trans-horizon correlations%
}{%
  Partial: $\widehat{\mathcal{M}}_O \subsetneq \mathcal{M}$.\\
  Global reconstruction requires\\
  compatible Oculara.%
}{%
  $\ker(\pi_O \circ F_O)$:\\
  all structure outside $H_O$\\
  and all sub-resolution\\
  variations within it.%
}
